\bookchapter{บทที่ 32: คู่มือฉุกเฉินด้านความปลอดภัยมือถือ}

\chapterintro{บทนี้คือ “แผนฉุกเฉิน” เราหวังว่าคุณจะไม่มีวันต้องใช้มันเลย แต่ถ้าเผลอเจอเหตุการณ์มือถือหาย ถูกขโมย หรือโดนหลอกลวง คู่มือเล่มนี้จะช่วยให้ความเสียหายน้อยที่สุด และเวลาภัยพิบัติอย่างแผ่นดินไหวหรือฝนตกหนักถล่มมา มือถือก็คือเครื่องมือรับการแจ้งเตือน ระบุตำแหน่งขอความช่วยเหลือ และติดต่อสื่อสาร แนะนำให้อ่านผ่านทั้งบทตั้งแต่ตอนนี้ แล้วแคปหน้าจอขั้นตอนสำคัญเก็บไว้ หรือสั่งพิมพ์วางไว้ที่บ้าน เพราะตอนเกิดเรื่องมือถืออาจไม่อยู่ในมือแล้ว ส่วนตอนภัยพิบัติแบตอาจหมดหรือสัญญาณอาจหาย}

\booksection{ฉากที่ 0: แผ่นดินไหว/สภาพอากาศสุดขั้ว - มือถือช่วยชีวิตคุณได้อย่างไร}

จุดเน้นคือความปลอดภัยของตัวคุณ ไม่ใช่บัญชีหรือเงิน ส่วนการตั้งค่าการแจ้งเตือนดูได้ใน บทที่ 13 “การแจ้งเตือนสภาพอากาศ ภัยธรรมชาติ และแผ่นดินไหว” ที่นี่จะพูดถึงแค่สิ่งที่ควรทำหลังเสียงแจ้งเตือนดังขึ้น

\textbf{เสียงแจ้งเตือนแผ่นดินไหวดังแล้ว: หลบก่อนแล้วค่อยอพยพ อย่ามัวแต่ถ่ายคลิป}\par

\begin{enumerate}
  \item หลบภัยใกล้ตัว: ก้มลงต่ำ หาที่กำบัง จับให้แน่น (ใต้โต๊ะ มุมผนังรับน้ำหนัก ข้างเฟอร์นิเจอร์ชิ้นใหญ่ที่แข็งแรง) ปกป้องศีรษะและคอ อยู่ห่างจากกระจก ผนังด้านนอก และโคมไฟห้อย
  \item อย่าพรวดพราดวิ่งลงตึก: แรงสั่นหลักมักใช้เวลาแค่สิบกว่าวินาทีถึงไม่กี่สิบวินาที บันไดอาจอันตรายกว่า และห้ามใช้ลิฟต์เด็ดขาด
  \item เมื่อแรงสั่นหยุดแล้วจึงอพยพอย่างเป็นระเบียบ: ใช้บันได ระวังแรงสั่นตามหลังและสิ่งของที่ร่วงหล่น ไปยังที่โล่ง ๆ และอยู่ห่างเสาไฟฟ้า ป้ายโฆษณา และอาคารที่เสียหาย
  \item ใช้มือถือแบบเอาอยู่: ส่งข้อความสั้น ๆ ให้ผู้ติดต่อฉุกเฉิน ("ฉันอยู่ที่XX ปลอดภัย/บาดเจ็บ ต้องการความช่วยเหลือ") ประหยัดแบตกว่าโทรยาว ถ้าส่ง SMS/WeChat ได้ให้ใช้ข้อความก่อน
  \item ถ้าไม่มีสัญญาณ: ดู บทที่ 7 เรื่อง SOS ผ่านดาวเทียม / ข้อความผ่านดาวเทียม (เฉพาะรุ่นที่รองรับ) หรือไปยังที่สูงหรือที่โล่งแล้วลองใหม่ ในบางสถานการณ์การโทรคุยผ่าน Bluetooth แบบไม่ใช้เน็ตก็ช่วยได้
  \item ปิดฟังก์ชันที่ไม่จำเป็น: ตอนฉุกเฉินเปิดโหมดเครื่องบินแล้วตั้งเวลาเปิด-ปิด ปิดการสแกน Bluetooth ลดความสว่างหน้าจอ จะช่วยให้แบตอึดขึ้นอีกหลายชั่วโมง เก็บไฟเหลือนิดเดียวไว้สำหรับส่งข่าวปลอดภัยและรับข้อมูลการช่วยเหลือ
\end{enumerate}

\begin{tcbwarn}{ฝนตกหนัก/พายุไต้ฝุ่น/น้ำท่วมฉับพลันตามไหล่เขา: ต้องลงมือตั้งแต่ขั้นแจ้งเตือน}
\begin{itemize}
  \item เห็นแจ้งเตือนระดับสีแดง: พยายามออกจากบ้านให้น้อยที่สุด อยู่ห่างพื้นที่ใต้ดิน สะพานลอด ลำน้ำ และทางลาดชัน
  \item ถ้าติดอยู่กลางแจ้ง: เดินขึ้นไปยังจุดสูงของอาคารที่แข็งแรง อย่าฝืนลุยน้ำ (น้ำท่วมอาจมีฝาปิดท่อหายหรือไฟรั่ว) เวลาโทร 119/110 ให้บอกจุดสังเกตและระดับน้ำให้ชัดเจน
  \item ก่อนออกเดินทาง: ใช้แอปสภาพอากาศและช่องแจ้งเตือนภัยของท้องถิ่นเช็กความเสี่ยง ถ้าขับรถไกลให้เลี่ยงเส้นทางเขาที่ลัดเลาะริมน้ำหรือริมหน้าผา
  \item ตกลงกันในครอบครัว: นัดหมายกันว่า "ถ้าติดต่อกันไม่ได้ ให้ใช้ตำแหน่งพิกัดล่าสุดจาก WeChat/SMS เป็นหลัก" เพื่อไม่ให้ทุกคนโทรรัวจนแบตหมดกันหมด
\end{itemize}

\end{tcbwarn}

\begin{tcbpractice}{4 สิ่งที่ควรทำเอาไว้ตั้งแต่ปกติ}
\begin{enumerate}
  \item เปิดการแจ้งเตือนแผ่นดินไหวของระบบ หรือแอปแจ้งเตือนที่เชื่อถือได้ และอนุญาตให้มันดังข้ามโหมดเงียบได้
  \item ปักหมุด "ผู้ติดต่อฉุกเฉิน" ไว้บนสุด 2-3 คน ให้แสดงได้ที่หน้าล็อกจอ (iOS/Android มีการ์ดฉุกเฉินทางการแพทย์หรือข้อความบนหน้าล็อกจอ)
  \item เตรียมพาวเวอร์แบงก์ไว้ที่บ้านและในกระเป๋า ก่อนเดินทางไกลชาร์จให้เต็มและพาสายชาร์จไปด้วย
  \item จำเบอร์ที่ "โทรได้โดยไม่ต้องเปิดสมุดรายชื่อ" สักหนึ่งเบอร์ (เบอร์มือถือคนในครอบครัว) แล้วจดใส่กระดาษเก็บไว้ในกระเป๋าสตางค์
\end{enumerate}

\end{tcbpractice}

\booksection{ฉากที่ 1: มือถือหาย/ถูกขโมย - 30 นาทีทอง}

30 นาทีแรกหลังมือถือหายคือ “ช่วงเวลาทอง” ถ้าจัดการถูกวิธีจะช่วยเงินทองและความเป็นส่วนตัวของคุณไว้ได้ ทำตามลำดับข้างล่างนี้ ห้ามข้ามขั้นตอนไหนเด็ดขาด:

\textbf{ขั้นตอนมาตรฐานภายใน 30 นาทีหลังมือถือหาย}\par

\begin{enumerate}
  \item ยืมมือถือคนอื่นโทรเข้าเบอร์ตัวเอง (นาทีที่ 1): อาจแค่ตกอยู่แถวนั้น ถ้าคนดี ๆ รับสายก็ได้คืน
  \item ล็อกมือถือจากระยะไกล (นาทีที่ 2-5): ใช้มือถือหรือคอมของคนอื่นเข้าสู่ระบบบริการค้นหาอุปกรณ์: iPhone: iCloud.com > Find My > เลือกอุปกรณ์ของคุณ > ทำเครื่องหมายว่าสูญหาย
  \item Android แต่ละแบรนด์: Huawei cloud.huawei.com / Xiaomi i.mi.com / OPPO cloud.oppo.com ฯลฯ
  \item แจ้งหาย SIM (นาทีที่ 5-10): โทรฝ่ายบริการลูกค้าของผู้ให้บริการ (China Mobile 10086/China Unicom 10010/China Telecom 10000) แจ้งว่ามือถือหายและขอระงับบริการ SIM ชั่วคราว ขั้นตอนนี้สำคัญที่สุด เพราะถ้ามิจฉาชีพได้ SIM ของคุณไป จะรับรหัสยืนยันและรีเซ็ตรหัสผ่านทุกอย่างของคุณได้
  \item ระงับบัญชีชำระเงินชั่วคราว (นาทีที่ 10-20): Alipay: โทร 95188 แจ้งหาย หรือล็อกอิน Alipay บนมือถือของคนอื่น > ของฉัน > การตั้งค่า > ศูนย์ความปลอดภัย > ระงับฉุกเฉิน
  \item WeChat: แจ้งตำรวจ 110 แล้วติดต่อฝ่ายบริการลูกค้า WeChat หรือไปที่ศูนย์ความปลอดภัย WeChat > ระงับบัญชี
  \item แอปธนาคาร: โทรเบอร์บริการลูกค้าที่พิมพ์อยู่ด้านหลังบัตร แล้วขอระงับการใช้งาน Internet Banking
  \item แจ้งคนใกล้ตัว (นาทีที่ 20-25): โพสต์ลง Moments บน WeChat หรือให้คนในครอบครัวช่วยส่งข้อความ แจ้งเพื่อนฝูงว่ามือถือหาย เพื่อกันมิจฉาชีพเอาตัวตนของคุณไปหลอกคนอื่น
  \item แจ้งความ (นาทีที่ 25-30): โทร 110 หรือไปที่สถานีตำรวจที่ใกล้ที่สุดเพื่อแจ้งความ และเก็บใบรับแจ้งความไว้ เพราะทั้งการเคลมประกันและข้อพิพาทกับธนาคารภายหลังต้องใช้ใบนี้
\end{enumerate}

\booksection{ฉากที่ 2: โดนหลอก - ช่วง 72 ชั่วโมงทองสำหรับการรับมือ}

เมื่อรู้ตัวว่าโดนหลอก สิ่งสำคัญที่สุดคือใจเย็น ๆ การตัดสินใจผิดพลาดตอนตื่นตระหนกจะทำให้เสียเพิ่มอีกรอบ ทำตามขั้นตอนต่อไปนี้:

\begin{enumerate}
  \item หยุดการโอนเงินทุกอย่างทันที: ไม่ว่ามิจฉาชีพจะพูดอะไร อย่าโอนเพิ่มแม้แต่บาทเดียว "โอนเงินประกันอีกก้อนแล้วเงินเดิมจะคืน" คือสคริปต์หลอกรอบสองแบบคลาสสิก
  \item เก็บหลักฐานทั้งหมด: อย่าลบประวัติแชท! อย่าลบประวัติการโทร! แคปหน้าจอบันทึกแชท บันทึกการโอนเงิน และข้อมูลบัญชีของอีกฝ่ายไว้ ทั้งหมดนี้คือหลักฐานสำคัญสำหรับการแจ้งความและการตามเงินคืน
  \item โทร 96110 ทันที: นี่คือสายด่วนต่อต้านมิจฉาชีพระดับประเทศ ที่ปรึกษาจะแนะนำวิธีหยุดโอนฉุกเฉิน ถ้าโอนไปไม่นาน (ปกติภายใน 24 ชั่วโมง) ธนาคารอาจระงับบัญชีฝ่ายตรงข้ามเพื่อดักเงินไว้ได้
  \item โทร 110 แจ้งความ: พามือถือ บัตรประชาชน และหลักฐานการโอนเงินไปที่สถานีตำรวจที่ใกล้ที่สุดเพื่อให้ปากคำ
  \item ติดต่อธนาคาร/แพลตฟอร์มชำระเงิน: แจ้งสถานการณ์ว่าโดนหลอก และขอระงับบัญชีที่เกี่ยวข้อง ถ้าเป็นเครดิตการ์ดถูกสวมรอยใช้ ให้ขอธนาคารเริ่มขั้นตอน chargeback
  \item เปลี่ยนรหัสผ่านทุกตัว: เปลี่ยนรหัสผ่านของ Apple ID/บัญชีแบรนด์ WeChat Alipay และแอปธนาคารจากอุปกรณ์อื่นที่ปลอดภัย แล้วเปิดการยืนยันตัวตนสองชั้นด้วย
\end{enumerate}

\booksection{ฉากที่ 3: มือถือติดไวรัส/โทรจัน - ทำอย่างไรดี}

ถ้าคุณพบว่ามือถือมีอาการแบบนี้ โฆษณาเด้งขึ้นมาเองโดยไม่ทราบสาเหตุ ติดตั้งแอปแปลก ๆ เองอัตโนมัติ ค่าโทรพุ่งขึ้นฉับพลัน หรือเน็ตถูกใช้ผิดปกติ ก็น่าจะติดมัลแวร์แล้ว:

\begin{itemize}
  \item ขั้นที่ 1: เข้าเซฟโหมด (Safe Mode) - บนมือถือ Android กดค้างปุ่มเปิด-ปิดเพื่อปิดเครื่อง แล้วกดค้างปุ่มเปิด-ปิดอีกครั้งเพื่อเปิดเครื่อง เมื่อเครื่องสั่นตอนบูต ให้กดค้างปุ่มลดเสียงไว้ไม่ปล่อยจนเข้าเซฟโหมด (มุมซ้ายล่างจะขึ้นคำว่า "Safe Mode") ในเซฟโหมดแอปบุคคลที่สามทุกตัวจะถูกปิดใช้งาน ช่วยให้ไล่หาตัวการได้ง่าย
  \item ขั้นที่ 2: ถอนการติดตั้งแอปที่น่าสงสัย - ในเซฟโหมดเข้า [การตั้งค่า] > [การจัดการแอป] จัดเรียงตามเวลา แล้วถอนแอปแปลก ๆ ที่เพิ่งติดตั้งล่าสุดและจำไม่ได้ว่ามาจากไหน
  \item ขั้นที่ 3: ติดตั้งแอปความปลอดภัยมาสแกน - ดาวน์โหลด Tencent Mobile Manager หรือ 360 Mobile Guard มาสแกนทั้งเครื่องหนึ่งรอบ
  \item ขั้นที่ 4: รีเซ็ตเครื่องกลับค่าโรงงาน (ทางเลือกสุดท้าย) - ถ้าวิธีข้างบนไม่ได้ผล ให้สำรองข้อมูลสำคัญก่อนแล้วค่อยรีเซ็ตระบบ ระวัง: การรีเซ็ตจะลบข้อมูลทั้งหมด ต้องสำรองรูปภาพ สมุดรายชื่อ และไฟล์สำคัญอื่น ๆ ให้เรียบร้อยก่อนเสมอ
  \item iPhone แทบไม่มีทางติดไวรัส ระบบ sandbox ของ iOS ทำให้มัลแวร์อยู่รอดยากมาก ถ้า iPhone ของคุณทำงานผิดปกติ ส่วนใหญ่มักเป็นบั๊กของแอปสักตัวหรือปัญหาการตั้งค่า ไม่ใช่ไวรัส การรีเซ็ตเครื่องกลับค่าโรงงานแก้ปัญหาได้ 99\%
\end{itemize}

\booksection{ฉากที่ 4: บัญชีถูกแฮ็ก - แย่งสิทธิ์ควบคุมคืนอย่างรวดเร็ว}

\begin{enumerate}
  \item เปลี่ยนรหัสผ่านทันที: ถ้ายังล็อกอินบัญชีได้อยู่ ให้เปลี่ยนรหัสผ่านเป็นอย่างแรก รหัสใหม่ต้องต่างจากรหัสเดิมสนิท
  \item ไล่อุปกรณ์อื่นออก: เข้าไปที่ "การจัดการอุปกรณ์ที่ล็อกอิน" ในการตั้งค่าความปลอดภัยของบัญชี แล้วดันอุปกรณ์แปลกหน้าออกจากระบบให้หมด
  \item ตรวจประวัติการล็อกอิน: ดูประวัติการล็อกอิน 7 วันล่าสุด ถ้ามีการล็อกอินจากตำแหน่งที่ไม่คุ้นเคย แปลว่าบัญชีถูกขโมยจริง
  \item เปิดการยืนยันตัวตนสองชั้น: ถ้ายังไม่เคยเปิด ให้เปิดเดี๋ยวนี้เลย นี่คือวิธีที่ได้ผลที่สุดในการกันบัญชีถูกขโมยซ้ำ
  \item ตรวจข้อมูลที่ผูกกับบัญชี: เช็กว่าเบอร์มือถือและอีเมลที่ผูกกับบัญชีถูกเปลี่ยนไปหรือเปล่า ถ้าถูกเปลี่ยน ให้ใช้ฟังก์ชันอุทธรณ์/ขอคืนบัญชีเพื่อกู้คืน
  \item ตรวจการตั้งค่าการชำระเงิน: ถ้าโดนขโมยคือ WeChat/Alipay ให้เช็กว่าการจ่ายแบบไม่ต้องใส่รหัสและการหักเงินอัตโนมัติถูกแก้ไขไหม และบัตรธนาคารถูกเพิ่มเข้าบัญชีแปลกหน้าหรือเปล่า
\end{enumerate}

\booksection{ฉากที่ 5: ได้รับข้อความกดเรียกค่าไถ่ - อย่ากลัว อย่าจ่าย}

แบบแผนการกดเรียกค่าไถ่ที่พบบ่อยมีสองแบบ: (1) “ฉันมีคลิป/รูปลับของคุณ ถ้าไม่จ่ายจะส่งให้ทุกคนในสมุดรายชื่อของคุณ” (2) “มือถือของคุณถูกแฮ็กเกอร์ควบคุมแล้ว ถ้าไม่จ่ายจะล็อกถาวร” เจอแบบนี้แล้วให้ทำดังนี้:

\begin{itemize}
  \item 99.9\% คือมิจฉาชีพ: พวกเขาไม่มีข้อมูลส่วนตัวของคุณอยู่ในมือเลย แค่ส่งข้อความกดเรียกค่าไถ่แบบเดียวกันเป็นหมื่นเป็นแสนฉบับ เดี๋ยวก็มีคนกลัวจนจ่าย
  \item อย่าตอบ อย่าจ่ายเงิน: การตอบกลับเท่ากับบอกอีกฝ่ายว่า "เบอร์นี้มีคนใช้อยู่" คุณจะถูกแท็กว่าเป็น "เป้าหมายที่กดเรียกค่าไถ่ได้" แล้วจะยิ่งโดนรบกวนต่อไปอีกเพียบ
  \item แคปหน้าจอเก็บหลักฐาน: แล้วบล็อกอีกฝ่ายทิ้ง
  \item เปลี่ยนรหัสผ่าน: ถ้ายังกังวลอยู่ ก็เปลี่ยนรหัสผ่านของบัญชีหลัก ๆ สักรอบเพื่อความอุ่นใจ
  \item ถ้ามีความเสี่ยงข้อมูลส่วนตัวรั่วไหลจริง ๆ: แจ้งความ ตำรวจมีประสบการณ์จัดการคดีแบบนี้
\end{itemize}

\booksection{ฉากที่ 6: มิจฉาชีพ AI เปลี่ยนใบหน้า/AI เลียนแบบเสียง - เห็นด้วยตาก็ไม่จริงอีกต่อไป}

ตั้งแต่ปี 2024 เป็นต้นมา เทคโนโลยี AI เปลี่ยนใบหน้าและ AI เลียนแบบเสียงทำให้มิจฉาชีพ “กลายเป็น” ใบหน้าและเสียงของครอบครัว หัวหน้า หรือเพื่อนของคุณในวิดีโอคอลได้แล้ว “แม่คะ ลูกมีเรื่อง รีบโอนเงินมาหน่อย” คนในคลิปหน้าตาและเสียงเหมือนลูกชายคุณเป๊ะ แต่จริง ๆ เป็นของปลอม นี่คือหนึ่งในมือของมิจฉาชีพที่กันยากที่สุดและอันตรายที่สุดในตอนนี้

\textbf{ขั้นที่ 1: หาจุดพร่อง - รอยตำหนิที่พบบ่อย 4 อย่างของงานปลอมจาก AI}\par

\begin{itemize}
  \item ความถี่การกะพริบตาผิดปกติ: คนจริงกะพริบตา 15-20 ครั้งต่อนาที แต่ใบหน้าปลอมจาก AI มักกะพริบน้อยเกินไปหรือจังหวะแข็ง ๆ ไม่เป็นธรรมชาติ
  \item ริมฝีปากกับเสียงไม่ตรงกัน: ดูดี ๆ จะเห็นว่ารูปปากกับเสียงที่ได้ยินหลุดกันอยู่เสี้ยววินาที
  \item ขอบใบหน้าเบลอ: กรอบหน้า แถวใบหู และเส้นขอบผมมักมี "ขอบหยัก ๆ" หรือกระพริบสั่นได้
  \item ทิศทางแสงเงาไม่สอดคล้อง: ไฮไลต์และเงาบนใบหน้าไม่ตรงกับทิศแสงของฉากหลัง หรือเวลาหันหน้าเฉียง อวัยวะส่วนหน้า "หมุนตาม" ไม่ถูกต้อง
\end{itemize}

\begin{tcbinfo}{ขั้นที่ 2: สามเทคนิคทดสอบว่าอีกฝ่ายจริงหรือปลอม}
\begin{itemize}
  \item ขอให้อีกฝ่ายทำท่าเฉพาะ: ให้เขาหันหัว เอามือปิดหน้า หรือใช้มือบังปากแล้วทำปากเปิด-ปิด ใบหน้าปลอมจาก AI จะ "หลุดบ่อง" ง่ายที่สุดตอนถูกบังหรือหันหัวเฉียงมุมใหญ่ ใบหน้าจะเบี้ยวหรือตามท่าไม่ทัน
  \item ถามคำถามที่มีแต่อีกฝ่ายรู้: แต่งคำถามขึ้นมา เช่น "คราวก่อนพาหมาบ้านเราไปรักษาที่คลินิกไหน?" ถ้าเป็นมิจฉาชีพ เขาจะตอบไม่ได้หรือพูดลุกลั่ม
  \item โทรกลับเบอร์ที่รู้จัก: วางสายวิดีโอ แล้วโทรกลับด้วยเบอร์ที่เก็บไว้ในสมุดรายชื่อของคุณโดยตรง มิจฉาชีพควบคุมเบอร์มือถือจริงของคนนั้นไม่ได้ โทรเมื่อไหร่ก็เผยทันที
\end{itemize}

\end{tcbinfo}

\begin{tcbwarn}{ขั้นที่ 3: ขั้นตอนแจ้งความหลังยืนยันว่าโดนหลอก}
\begin{enumerate}
  \item หยุดการโอนเงินทันที: ยังไงก็ "ด่วน" แค่ไหนก็ให้หยุดก่อน ใครก็ตามที่รีบให้ "โอนเร็วเข้า อย่าวางสาย" แทบทุกคนคือมิจฉาชีพ
  \item เก็บหลักฐานทั้งหมด: อัดหน้าจอ แคปภาพหน้าจอวิดีโอคอล บันทึกแชท บัญชีของอีกฝ่าย และหลักฐานการโอนเงิน ทั้งหมดนี้คือวัสดุพยานสำคัญสำหรับการตรวจพิสูจน์ทางเทคนิค
  \item โทรสายด่วนต่อต้านมิจฉาชีพ 96110: แจ้งว่าเป็นมิจฉาชีพประเภท "AI เปลี่ยนหน้า/เลียนเสียง" ทางสายด่วนจะเริ่มขั้นตอนหยุดโอนฉุกเฉินให้
  \item โทร 110 แจ้งความ: พาหลักฐานครบไปที่สถานีตำรวจเพื่อให้ปากคำ โดยระบุชัดว่าฝ่ายนั้นใช้เทคโนโลยีปลอมแปลงด้วย AI เพื่อให้ตำรวจเปิดคดีและจัดชนิดคดีได้ถูกต้อง
  \item แจ้งคนใกล้ตัว: มิจฉาชีพอาจเอาวิดีโอปลอมของคุณไปหลอกคนอื่น ลองโพสต์เตือนลง Moments หรือกลุ่มครอบครัวไว้ก่อน
\end{enumerate}

\end{tcbwarn}

\booksection{ฉากที่ 7: การโจมตี SIM swap (การแย่งชิงซิม) - เบอร์ของคุณถูก “ขโมย” ไปแล้ว}

นี่คือการโจมตีที่หลายคนไม่เคยได้ยิน แต่สร้างความเสียหายร้ายแรงมาก SIM swap (การแย่งชิงซิมการ์ด) คือการที่ผู้โจมตีปลอมข้อมูลตัวตนของคุณ แล้วแอบอ้างเป็นคุณไปที่ศูนย์บริการเพื่อทำซิมใหม่แทนเบอร์เดิม พอทำสำเร็จ สัญญาณมือถือของคุณจะหายไปพรวดเดียว ส่วนมือถือของผู้โจมตีที่เสียบซิมใบใหม่นั้นจะรับรหัสยืนยันทาง SMS ทั้งหมดของคุณได้ ทั้ง Alipay ธนาคาร WeChat ประตูบานใหญ่ทุกบานเปิดโล่งหมด

\begin{tcbinfo}{จะรู้ได้อย่างไร? สัญญาณเตือน 3 อย่างของ SIM swap}
\begin{itemize}
  \item มือถือขึ้น "ไม่มีสัญญาณ"/"ไม่พบซิมการ์ด" ทั้งที่คุณไม่ได้แตะซิมเลย และรีสตาร์ตก็ไม่หาย มักเป็นช่วงเวลาที่ซิมถูกแย่งไปพอดี
  \item ได้รับ SMS จากผู้ให้บริการว่า "ทำซิมใหม่สำเร็จ" หรือ "ทำรายการสำเร็จ" ทั้งที่คุณไม่ได้ทำเองเลย
  \item ทันทีทันใดรับรหัสยืนยันไม่ได้ แถม SMS อื่น ๆ ก็หายหมด เพราะข้อความถูกดักไปที่ซิมของผู้โจมตีแล้ว
\end{itemize}

\end{tcbinfo}

\begin{tcbwarn}{มาตรการป้องกัน - ใส่กุญแจให้ซิมการ์ดด้วย}
\begin{itemize}
  \item ตั้งรหัส PIN ให้ซิมการ์ด: เส้นทาง [การตั้งค่า] > [โทรศัพท์] > [รหัส PIN ซิมการ์ด] เปิดใช้แล้ว ทุกครั้งที่รีสตาร์ตเครื่องหรือเปลี่ยนซิม ต้องกรอก PIN ก่อนถึงจะใช้ซิมนี้ได้ ผู้โจมตีได้ซิมไปก็ใช้ไม่ได้
  \item เปิดการยืนยันสองชั้นกับผู้ให้บริการ: ล็อกอินแอป China Mobile/China Unicom/China Telecom แล้วหา "ศูนย์ความปลอดภัย" หรือ "การยืนยันทำรายการ" เปิด SMS แจ้งเตือนเมื่อมีคนอื่นทำซิมแทน และยืนยันรหัสผ่านรอบสองสำหรับการทำรายการ อย่างนี้ใครไปทำซิมแทนคุณ พนักงานจะขอตรวจสอบตัวตนเพิ่ม
  \item บัญชีสำคัญอย่าพึ่ง SMS อย่างเดียว: บัญชีมูลค่าสูงอย่างธนาคารหรือ Apple ID ควรใช้แอปยืนยันตัวตน (Google Authenticator, Microsoft Authenticator) หรือฮาร์ดแวร์คีย์เป็นหลัก ไม่ใช่รหัสยืนยันทาง SMS
  \item ใส่ลายน้ำบนสำเนาบัตรประชาชน: สำเนาที่ใช้ทำรายการให้เขียนกำกับว่า "ใช้สำหรับรายการ XX เท่านั้น" เพราะผู้โจมตีมักอาศัยข้อมูลบัตรประชาชนที่รั่วไหลไปปลอมเป็นคุณ
\end{itemize}

\end{tcbwarn}

\textbf{การรับมือฉุกเฉินหลังซิมถูกแย่ง - ความเร็วคือทุกอย่าง}\par

\begin{enumerate}
  \item โทรฝ่ายบริการลูกค้าของผู้ให้บริการเพื่อแจ้งหายทันที: China Mobile 10086 / China Unicom 10010 / China Telecom 10000 ขอระงับเบอร์และระงับบริการทำซิมใหม่ทันที เพื่อกันไม่ให้ผู้โจมตีใช้เบอร์ของคุณต่อ
  \item ล็อกอินบัญชีสำคัญจากอุปกรณ์อื่น: ใช้คอมหรือมือถือของคนในครอบครัว เปลี่ยนรหัสผ่านของธนาคาร Alipay WeChat Apple ID ทันที แล้วไล่อุปกรณ์แปลกหน้าออกด้วย
  \item ระงับบัญชีชำระเงิน: Alipay โทร 95188 แจ้งหาย, WeChat ระงับผ่านศูนย์ความปลอดภัย, ธนาคารให้ฝ่ายบริการลูกค้าระงับ Internet Banking
  \item พาบัตรประชาชนไปศูนย์บริการ: ไปด้วยตัวเองพร้อมบัตรประชาชนที่ศูนย์บริการหลักของผู้ให้บริการเพื่อทำซิมใหม่ แล้วกลับมาใช้เบอร์เดิมได้ตามปกติ
  \item แจ้งความ: SIM swap ถือเป็นความผิดฐานละเมิดข้อมูลส่วนบุคคลแล้ว พาหลักฐานให้ครบ (บันทึกเวลาที่เบอร์ถูกระงับ, ความเสียหายทางการเงิน, บันทึกการทำซิมใหม่) ไปแจ้งความ
  \item ตรวจสอบทั้งระบบ: หลังกลับมาใช้เบอร์ได้ ไล่เช็กประวัติการล็อกอินและรายการเงินช่วงล่าสุดทีละรายการ เพื่อยืนยันว่าไม่มีร่องรอยถูกเอาตัวตนไปใช้
\end{enumerate}

\begin{tcbwarn}{กันไว้ดีกว่าแก้ - 5 สิ่งที่ทำได้ในชีวิตประจำวัน}
\begin{enumerate}
  \item รหัสล็อกจอ + การยืนยันตัวตนสองชั้น เปิดทั้งสองอย่าง
  \item ยืนยันว่า "ค้นหามือถือของฉัน" เปิดอยู่
  \item ตั้ง PIN ให้ซิมการ์ด ([การตั้งค่า] > [โทรศัพท์] > [รหัส PIN ซิมการ์ด]) มือถือหายแล้วคนอื่นจะใช้ซิมของคุณไม่ได้
  \item ไม่ติดตั้งแอปที่ไม่รู้ที่มา ไม่คลิกลิงก์ของคนแปลกหน้า
  \item สำรองข้อมูลสำคัญสม่ำเสมอ (รูปภาพ สมุดรายชื่อ บันทึกแชท)
\end{enumerate}

\end{tcbwarn}

\begin{tcbpractice}{ลงมือทำ: ทำ "บัตรติดต่อฉุกเฉิน" ของคุณเอง}
\begin{enumerate}
  \item หากระดาษมาแผ่นหนึ่ง จดเบอร์เหล่านี้: ฝ่ายบริการลูกค้าผู้ให้บริการ (China Mobile 10086/China Unicom 10010/China Telecom 10000) ฝ่ายบริการลูกค้าธนาคาร สายด่วนต่อต้านมิจฉาชีพ 96110 แจ้งเหตุด่วน 110 และเบอร์คนในครอบครัว
  \item จดอีเมลที่ใช้สมัครและช่องทางติดต่อสำรองของ Apple ID/บัญชีแบรนด์ของคุณ
  \item เก็บกระดาษแผ่นนี้ไว้ในกระเป๋าสตางค์หรือแปะไว้ที่ตู้เย็น อย่าเก็บไว้ในมือถือ เพราะตอนเกิดเรื่องมือถืออาจไม่อยู่กับตัว
  \item เช็กเดี๋ยวนี้เลย: "ค้นหามือถือของฉัน" เปิดแล้วหรือยัง และ PIN ของซิมการ์ดตั้งไว้แล้วหรือยัง
\end{enumerate}

\end{tcbpractice}

\begin{tcbsummary}{สรุปบทนี้}
\begin{itemize}
  \item 30 นาทีทองหลังมือถือหาย: ยืมมือถือโทรเข้าเบอร์ตัวเอง ล็อกระยะไกล แจ้งหายซิม ระงับบัญชีชำระเงิน แจ้งคนใกล้ตัว แจ้งความ
  \item 72 ชั่วโมงหลังโดนหลอก: หยุดโอนเงิน เก็บหลักฐาน โทร 96110 แจ้งความ ติดต่อธนาคาร เปลี่ยนรหัสผ่าน
  \item มือถือติดไวรัส: เซฟโหมด ถอนแอปน่าสงสัย สแกนความปลอดภัย รีเซ็ตค่าโรงงาน (ทางเลือกสุดท้าย)
  \item บัญชีถูกขโมย: เปลี่ยนรหัสผ่าน ไล่อุปกรณ์ ตรวจประวัติล็อกอิน เปิดยืนยันสองชั้น ตรวจข้อมูลผูกบัญชี ตรวจการชำระเงิน
  \item ข้อความกดเรียกค่าไถ่ 99.9\% เป็นมิจฉาชีพ ไม่ตอบ ไม่จ่ายเงิน แคปหลักฐานแล้วบล็อก
  \item มิจฉาชีพ AI เปลี่ยนหน้า/เลียนเสียง: จับผิดจากการกะพริบตา/รูปปาก/ขอบใบหน้า/แสงเงา ขอให้อีกฝ่ายหันหัวปิดหน้า ถามคำถามลับ และโทรกลับเบอร์ที่รู้จัก
  \item การแย่งชิง SIM swap: เมื่อมือถือขึ้น "ไม่มีสัญญาณ" ฉับพลันให้ตื่นตัวทันที ตั้ง PIN ซิม เปิดยืนยันสองชั้นกับผู้ให้บริการ และเปลี่ยนบัญชีสำคัญมาใช้แอปยืนยันตัวตน
  \item การป้องกันประจำวัน 5 อย่าง: ล็อกจอ+2FA ค้นหามือถือ PIN ซิม ไม่ติดตั้งแอปไม่รู้ที่มา และสำรองข้อมูลสม่ำเสมอ
\end{itemize}

\end{tcbsummary}
